\documentclass[12pt]{article}
\usepackage{amsmath, amssymb, amsthm, epsfig}
\usepackage{listings}
\usepackage{fontspec}
\setmonofont{DejaVu Sans Mono}
\usepackage{tikz}
\usepackage[shortlabels]{enumitem}
\usetikzlibrary{math,decorations.pathreplacing, shapes.geometric}
\usepackage{quiver}
\usepackage{mathtools}
\usepackage{float}
\usepackage{subcaption}

\usepackage[T1]{fontenc}
\usepackage[utf8]{inputenc}
\usepackage[fixed]{fontawesome5}
\usepackage{pifont}
\usepackage{bbding}

\theoremstyle{definition}
\usepackage{enumitem}
\newenvironment{Proof}{\noindent {\sc Proof.}}{$\Box$ \vspace{2 ex}}
\newenvironment{Solution}{\noindent {\sc Solution.}}{$\Box$ \vspace{2 ex}}


\usepackage{color}
\definecolor{keywordcolor}{rgb}{0.7, 0.1, 0.1}   % red
\definecolor{tacticcolor}{rgb}{0.0, 0.1, 0.6}    % blue
\definecolor{commentcolor}{rgb}{0.4, 0.4, 0.4}   % grey
\definecolor{symbolcolor}{rgb}{0.0, 0.1, 0.6}    % blue
\definecolor{sortcolor}{rgb}{0.1, 0.5, 0.1}      % green
\definecolor{attributecolor}{rgb}{0.7, 0.1, 0.1} % red

\def\lstlanguagefiles{../lstlean.tex}

\oddsidemargin-1mm
\evensidemargin-0mm
\textwidth6.5in
\topmargin-15mm
\textheight8.75in
\footskip27pt

\usepackage[table]{xcolor}
\definecolor{truecolor}{RGB}{198,239,206} % Light green
\definecolor{falsecolor}{RGB}{255,199,206} % Light red
\newcommand{\true}{\cellcolor{truecolor}T}
\newcommand{\false}{\cellcolor{falsecolor}F}
\newcommand{\dom}{\mathrm{dom}}
\begin{document}

% homework problems are from the following textbook:
% Larry J. Gerstein, Introduction to Mathematical Structures and Proofs, Undergraduate Texts in Mathematics (https://link.springer.com/book/10.1007/978-1-4614-4265-3)
% Read Sections 3.1--3.3 (Gerstein's book).
% §3.1) 4, 6, 7, 11
% §3.2) 4, 5, 10, 13
% §3.3) 7, 8, 10, 14
% NOTE: We are using the definition that
% A function $f : A \to B$ is a relation $f \subseteq A \times B$ such that $\forall (x \in A), \exists! (y \in B), (x, y) \in f$.
% Thus, two functions are equal if their graphs are equal as sets.
\noindent \textit{\textbf{Math 300, Winter 2026}} \hspace{1.3cm} \textit{\textbf{HOMEWORK 4}} \hspace{1.3cm} 
\textit{\textbf{Grant Yang}
} 

\vspace{0.5cm}

\begin{Solution}
{\bf (\S3.1, Problem 4)}\\
\begin{enumerate}[(a)]
    \item Let $f(x)=1$ and $g(x)=\frac{x-5}{x-5}$. Then we have $f:\mathbb{R}\to \mathbb R$ and $g : \mathbb{R} \setminus \{5\}\to \mathbb R$, so $f \ne g$ because $\dom\,f \ne \dom\,g$.
    \item Let $f(x)=\sqrt x$ and $g(x)=\sqrt{|x|}$. Then $f : \mathbb R_{\ge 0} \to \mathbb R$ and $g : \mathbb{R} \to \mathbb{R}$, so $f \ne g$ because $\dom\,f \ne \dom\,g$.
    \item Let $f(x) = |x|$ and $g(x) = \sqrt{x^2}$. Then $f : \mathbb{R} \to \mathbb{R}$ and $g : \mathbb{R} \to \mathbb{R}$, and $f = g$ because $\forall(x\in\mathbb R), |x|=\sqrt{x^2}$.
    \item Let $f(x)=x^2-x-6$ and $g(x)=(x-4)(x+3)+6$. Then $f : \mathbb{R} \to \mathbb R$ and $g : \mathbb{R} \to \mathbb{R}$, and $f = g$ because 
    \begin{align*}
        \forall (x \in \mathbb R),\; x^2-x-6&=(x-4)(x+3)+6
        \\ &= x^2 -4x+3x-12+6
        \\ &= x^2 - x - 6. \qquad\square
    \end{align*}
    \item Let $f(x)=x^2$ and $g(x)=\begin{cases}x^2 & \text{if } x \text{ is rational} \\ 0 & \text{if } x \text{ is irrational}\end{cases}$. We have $f : \mathbb{R} \to \mathbb{R}$ and $g : \mathbb{R} \to \mathbb{R}$, but $f \ne g$ because for $x=\sqrt{2}$, we have $f(x)=2$ but $g(x)=0$.
\end{enumerate}
\end{Solution}

\vspace{0.5cm}

\begin{Solution}
{\bf (\S3.1, Problem 6)}\\
Given a function $f : A \to B$, define a relation $\sim$ by $a_1 \sim a_2 \iff f(a_1) = f(a_2)$.
\begin{enumerate}[(a)]
    \item This relation is an equivalence relation. We have transitivity, that $a_1 \sim a_2 \land a_2 \sim a_3 \implies a_1\sim a_3$. By the definition of our relation, we have from the hypothesis that $f(a_1) = f(a_2)$ and that $f(a_2) = f(a_3)$. Since $=$ is itself an equivalence relation and is thus transitive, we have $f(a_1)=f(a_3)$, so $a_1 \sim a_3$. Additionally, by reflexivity of $=$, we have that $f(a_1) = f(a_1)$, implying reflexivity for $\sim$ too (i.e. $a_1 \sim a_1$). Finally, we have symmetry, that $a_1 \sim a_2 \implies a_2 \sim a_1$, again from the symmetry of $=$. From the hypothesis we deduce $f(a_1)=f(a_2)$ and thus $f(a_2) = f(a_1)$ and finally that $a_2 \sim a_1$.
    \item Let $A$ be the set of solid-color cars and let $B$ be the set of all colors, with $f(x)$ being the color of $x$. Then $\sim$ induces a partition by color, where the equivalence classes are of cars that all have the same color.
    \item Let $A=B=\mathbb R$ and $f(x)=x^2$. This partitions the reals by magnitude/absolute value, where the equivalence class of $x$ is $[x]=\{x,-x\}$ (and $[0]$ is just the singleton $\{0\}$).
    \item Let $A=B=\mathbb R$ with $f(x)=|x|$. This induces the same partition as $f(x)=x^2$ in part c, where each equivalence class consists of $\{x, -x\}$.
    \item Let $A=\mathbb R \times \mathbb R$ and $B=\mathcal P(\mathbb R)$ with $f(x,y)=\{x,y\}$. This partition's equivalence classes are of the form $[(x,y)]=\{(x,y),(y,x)\}$, equating each point with its reflection across the $y=x$ line.
    \item Let $A=\mathbb R \times \mathbb R$ and $B = \mathbb R$ with $f(x,y) = x+y$. This partitions the $xy$-plane into equivalence classes that are diagonal lines of the form $y = -x+b$.
    \item Let $A=\mathbb R \times \mathbb R$ and $B = \mathbb R$ with $f(x,y)=x^2+y^2$. This partitions the $xy$-plane into equivalence classes that are circles centered at $(0,0)$.
\end{enumerate}
\end{Solution}

\vspace{1cm}

\begin{Solution}
{\bf (\S3.1, Problem 7)}\\
Define the characteristic function $\chi_A:S \to \{0,1\}$ for some $A\subseteq S$ to be $\chi_A(x)=\begin{cases}1 & x \in A \\ 0 & x \not \in A\end{cases}$. 
\begin{enumerate}[(a)]
    \item We have $\chi_\varnothing(x)=0$ and $\chi_S(x)=1$.
    \item Let $A'$ be the complement of $A$ in $S$. Then $\chi_A(x)=1-\chi_{A'}(x)$.
    \item $A=B\iff \chi_A=\chi_B$. Consider the forward direction. We can see that $\dom\,\chi_A=\dom\,\chi_B=S$. Let $x\in S$, and first let us consider the case that $x \in A$. In this case, $\chi_A(x)=1$ by definition, and from the definition of set equality, we get that $x \in B$ and $\chi_B(x)=1$. Next, consider the remaining case that $x \not\in A$. Then $\chi_A(x)=0$, and since $x\not\in B$, $\chi_B(x)=0$. Thus, we have that $\forall(x\in S),\chi_A(x)=\chi_B(x)$ and so $\chi_A = \chi_B$.

    Now consider the backwards direction. Let $x \in A$. Then by definition of $\chi_A$ we have $\chi_A(x)=\chi_B(x)=1$, so by the definition of $\chi_B$ we have $x \in B$ and thus $A \subseteq B$. Let $x \in B$. Then again we have $\chi_A(x)=\chi_B(x)=1$, and so $x \in B$ and $B \subseteq A$. Thus, $A=B$.

    \item \[\forall(x\in S), \chi_{A \cap B}(x)= \chi_A(x) \cdot \chi_B(x)\]
    Let $x \in S$. If $x \in A \land x \in B$, we have $\chi_A(x) = 1$ and $\chi_B(x)=1$ so $\chi_A(x)\cdot\chi_B(x)=1$. Additionally, $x \in A \cap B$, so $\chi_{A\cap B} (x)=1$ as desired. Now consider the remaining case that $\lnot(x \in A\land x\in B)$. This implies that either $x\not\in A$ or $x \not\in B$, so $x \not\in A \cap B$ and $\chi_{A\cap B}(x)=0$. Additionally, since we either have $\chi_A(x)=0$ or $\chi_B(x)=0$, we must have  $\chi_A(x)\cdot\chi_B(x)=0$ as desired. Thus, $\chi_{A\cap B}(x)=\chi_A(x)\cdot\chi_B(x)$.

    \[\forall(x\in S), \; \chi_{A\cup B}(x) = \chi_A(x)+\chi_B(x)-\chi_A(x)\chi_B(x).\]

    Let $x \in S$. Consider the case that $x \in A \lor x \in B$, and within this case consider $x \in A \land x \in B$. Then $\chi_{A\cup B}(x) = \chi_A(x)=\chi_B(x)= 1$, and we have $1 = 1 + 1 - 1$ as desired. Now consider the (sub)case that $\lnot(x \in A \land x\in B)$. Then we either have that $x \not\in A$ or $x \not\in B$ (but not both, as we still have that $x \in A \lor x \in B$), so $\chi_A(x)+\chi_B(x)=1=\chi_{A\cup B}(x)$ and $\chi_A(x)\cdot\chi_B(x)=0$, and so $1 = 1 - 0$ as desired. Finally, consider the case that $\lnot(x \in A \lor x\in B)$, so that $x\not\in A$ and $x \not\in B$. Then, $x \not\in A \cup B$, so $\chi_{A\cup B}(x)=\chi_A(x)=\chi_B(x)=0$, so we have $0 = 0+0-0$ as desired.
    Yay principle of inclusion/exclusion!
\end{enumerate}
\end{Solution}

\vspace{0.5cm}

\begin{Solution}
{\bf (\S3.1, Problem 11)}\\
Let $B^A$ denote the set of functions $A\to B$.
\begin{enumerate}[(a)]
    \item $\{1,2\}^{\{1,2\}} = \{\{(1,1),(2,1)\}, \{(1,1),(2,2)\},\{(1,2),(2,1)\},\{(1,2),(2,2)\}\}$.
    \item If $A \subseteq B$, then $A^C \subseteq B^C$. Let $f \in A^C$ and let $c \in C$. Then $f(c) \in A$, and by the definition of subset, $f(c) \in B$. Thus, in addition to considering $f$ to be $f : C \to A$, we can also consider it to be $f : C \to B$, and thus $f \in B^C$ and $A^C \subseteq B^C$.
    \item The sets $\{1,2\}^{\{1,2\}}$ and $\{1,2\}^{\{1,2,3\}}$ are disjoint. Let $g \in \{1,2\}^{\{1,2,3\}}$ be a function $g : \{1,2,3\}\to\{1,2\}$. By the definition of function, there must exist some (unique) $y\in \{1,2\}$ such that $(3,y)\in g$. However, since $3\not\in \{1,2\}$, we have that $g \not\subseteq \{1,2\}\times\{1,2\}$, so $g$ is not a relation on $\{1,2\}\times\{1,2\}$ and thus is not a function $\{1,2\}\to\{1,2\}$. We have shown that $\forall(g\in \{1,2\}^{\{1,2,3\}}),g\not\in \{1,2\}^{\{1,2\}}$, so $\{1,2\}^{\{1,2\}} \cap \{1,2\}^{\{1,2,3\}}=\varnothing$.
    \item If $B \ne C$, then $A^B \cap A^C = \varnothing$. Assume for contradiction that there exists some $f \in A^B \cap A^C$. Because $f$ is simultaneously a function $B \to A$ and a function $C \to A$, we must have that $\dom\,f = B$ and $\dom\,f = C$, but this contradicts the fact that $B \ne C$. Thus, no such $f \in A^B \cap A^C$ exists, and $A^B \cap A^C = \varnothing$.
    \item If $A$ is a nonempty set, then $\varnothing^A = \varnothing$. Let $x\in A$, and assume for contradiction that we have $f \in \varnothing^A$. Then by the definition of a function we must have that $f(x)=n$ for some unique $n \in \varnothing$. However, there does not exist any elements of the empty set, so our assumption was false and there does not exist any $f \in \varnothing^A$. Thus, $\varnothing^A = \varnothing$.
    \item $B^\varnothing = \{\varnothing\}$. Let $f \in B^\varnothing$ be any function $f : \varnothing \to B$. By the definition of a function, we have that $f \subseteq \varnothing \times B$, but since $\varnothing \times B = \varnothing$, we get that in fact $f = \varnothing$ (as $\varnothing \subseteq f$ is true for any set). Thus, $B^\varnothing \subseteq \{\varnothing\}$, but we still need to show that $\{\varnothing\} \subseteq B^\varnothing$, i.e. $\varnothing$ is a valid function $\varnothing : \varnothing \to B$. We have that $\varnothing \subseteq B\times \varnothing$, and we also have that $\forall (n\in \varnothing),\exists!(b\in B), \varnothing(b)=n$. This statement is trivially true as there does not exist any $n\in \varnothing$, so the $\forall$ is automatically satisfied. Thus, $\varnothing$ is a function and $B^\varnothing = \{\varnothing\}$.
\end{enumerate}
\end{Solution}

\vspace{1cm}

\begin{Solution}
{\bf (\S3.2, Problem 4)}\\
Let $S$ and $T$ be sets, with $|S|=3$ and $|T|=2$.
\begin{enumerate}[(a)]
    \item The number of functions $f : S \to T$: each function must decide a value of $T$ for each $S$, so there are $|T|^{|S|}= 8$ functions.
    \item The number of injections $f : S \to T$: By the pigeonhole principle, it is impossible to assign a unique output to three inputs from 2 possible values. There are 0 injections.
    \item The number of injections $f : T \to S$: We have 3 choices for the first element of $T$, and then 2 remaining choices, so there are $3\cdot 2 = 6$ injections.
    \item The number of surjections $f : S\to T$: Of the 8 functions from part (a), only the 2 constant functions fail surjectivity. So there are $8-2=6$ surjections.
    \item The number of surjections $f : T\to S$: There are 0 surjections as it is impossible to cover $|S|=3$ elements using only $|T|=2$ values.
    \item Let $|S|=m$ and $|T|=n$. Then the number of functions $S\to T$ is $|T|^{|S|} = n^m$ and the number of injections is $\frac{|T|!}{(|T|-|S|)!} = \frac{n!}{(n-m)!}$ (assuming $m \le n$, otherwise 0). The number of surjections is related to THE STIRLING NUMBERS OF THE SECOND KIND and is given by $n!\cdot\genfrac{\{}{\}}{0pt}{}{m}{n} = \sum_{k=0}^n (-1)^{n-k}\binom{n}{k} k^m$.
\end{enumerate}
\end{Solution}

\vspace{0.5cm}

\begin{Solution}
{\bf (\S3.2, Problem 5)}\\
A function $f : A \to B$ is a subset $f \subseteq A \times B$, and so it is a relation as $f \subseteq (A \cup B) \times (A \cup B)$. We thus can consider the inverse relation $f^{-1} = \{(b,a) \mid (a,b) \in f\}$.
\begin{enumerate}[(a)]
    \item Consider $g,h : \{1,2,3\} \to \{-2,5,6\}$ given by $g = \{(1,5),(2,-2),(3,6)\}$ and $h=\{(3,5),(2,-2),(1,5)\}$. The inverses are given by $g^{-1} = \{(5,1),(-2,2),(6, 3)\}$ and $h^{-1}=\{(5,3),(-2,2),(5,1)\}$.
    \item If a function $f : A\to B$ is not injective, then the inverse $f^{-1}$ is not a function. By the definition of non-injective, we know that there exist $a_1,a_2 \in A$ such that $f(a_1) = f(a_2)$ but $a_1 \ne a_2$. This means that $(a_1, f(a_1)) \in f$ and $(a_2, f(a_2)) \in f$, so $(f(a_1),a_1)\in f^{-1}$ and $(f(a_2),a_2) \in f^{-1}$. However, since $f(a_1) = f(a_2)$ this implies that $f^{-1}$ is not a function: the $y \in A$ such that $f^{-1}(f(a_1)) = y$ is not unique: we can have $y=a_1$ or $y=a_2$, and $a_1 \ne a_2$.
    \item If $f$ is injective but not surjective, then $f^{-1} : f(A) \to A$ is a function with $f(A) \subsetneq B$. Let $(y,x_1) \in f^{-1}$ be any element of $f^{-1}$. Assume for the sake of contradiction that there exists some $x_2 \in A$ with $x_1 \ne x_2$ such that $(y, x_2) \in f^{-1}$. By the definition of inverse relation, this would imply that $(x_1, y) \in f$ and $(x_2, y) \in f$ so $f(x_2)=f(x_1)=y$. By the injectivity of $f$, we get that $x_1 = x_2$. However, this contradicts our assumption, so we know that $f^{-1}$ must be a function as the assignment of $x_1 = f(y)$ is unique for all $y \in f(A)$. 

    Now we show that we have a proper subset $f(A) \subsetneq B$. Let $y\in f(A)$ be any element, and by the fact that $f(A)$ is the domain, we must have $(y, x) \in f^{-1}$ for some (unique) $x \in A$. By the definition of inverse relation, we know that $(x,y) \in f$, and so $y \in B$ and $f(A) \subseteq B$. By the failure of surjectivity, we know that there exists some $b \in B$ such that $\forall (a \in A), f(a) \ne b$. This means that $b \not\in f(A)$: if it were, we could apply the definition of $f(A)$ as the domain of the inverse relation and get $(x, b) \in f$ for some $x \in A$, but this contradicts the defining property of $b$. Thus, $b \in B$ but $b \not\in f(A)$, so $B \ne f(A)$, and we have $f(A) \subsetneq B$ as desired.
\end{enumerate}
\end{Solution}

\vspace{0.5cm}

\begin{Solution}
{\bf (\S3.2, Problem 10)}\\
Let $\Pi = \mathbb R^2$ be the $xy$-plane, let $O=(0,0)$ be the origin, and let $\Pi^\ast = \Pi - \{O\}$ be the punctured plane. Let $L$ be the set of all lines through $O$. Define $f : \Pi^\ast\to L$ to be the function taking all points $P \in \Pi^\ast$ to $f(P)$, the line containing the segment $OP$. This is not a bijection because it is not injective: $f((0,1))=f((0, 2))$ but $(0,1)\ne(0,2)$. However, we can define $S=\{(\cos t, \sin t) \mid t \in [0, \pi)\} \subset \Pi^\ast$, and $f|_S : S \to L$ is a bijection. 

I love $\mathbb {RP}^1$. I think choosing $S = \{(x,1) \mid x\in \mathbb R\} \cup \{(1,0)\}$ also works and is funnier.
\end{Solution}

\vspace{0.5cm}

\begin{Solution}
{\bf (\S3.2, Problem 13)}\\
Let $\mathbb{Z}_2^n$ denote the space of $n$-bit long binary codewords. Define the Hamming distance $d : \mathbb{Z}_2^n \times \mathbb Z_2^n \to \mathbb Z$ to be
\[ d(a_1a_2\dots a_n, b_1b_2\dots b_n) = \sum_{i=1}^n |a_i-b_i|. \]
\begin{enumerate}[(a)]
    \item $d(s_1,s_1) \ge 0$. By the definition of Hamming distance, $d(s_1,s_2)$ is always a sum of non-negative integers (absolute value), and thus $d(s_1,s_2)$ is non-negative.
    \item $d(s_1,s_2) = 0 \iff s_1 = s_2$. Consider the forward direction. Assume for contradiction that $s_1 \ne s_2$. Then there must be a position $i$ at which $s_1$ and $s_2$ differ, and thus the Hamming distance includes either the term $|1 - 0| =1$ or $|0 - 1| = 1$. Thus, $d(s_1,s_2)\ge 1$, but this contradicts our hypothesis that $d(s_1,s_2)=0$. Therefore, we have $s_1 = s_2$. Now consider the backwards direction: because $s_1 = s_2$, we have that $\forall i, a_i = b_i$ in the Hamming sum, so we get $d(s_1,s_2)=\sum_{i=1}^n |a_i -a_i| = \sum_{i=1}^n 0 = 0.$
    \item $d(s_1,s_2)=d(s_2,s_1)$. Expanding definitions:
    \begin{align*}
        d(a_1a_2\dots a_n, b_1b_2\dots b_n) &= \sum_{i=1}^n |a_i-b_i| \\
         &= \sum_{i=1}^n |b_i-a_i|
         \\ &=d(b_1b_2\dots b_n, a_1a_2\dots a_n).
    \end{align*}
    Thus, $d(s_1, s_2)=d(s_2,s_1)$.
    \item $\forall(s_1,s_2,s_3\in \mathbb{Z}_2^n),\;d(s_1,s_3)\le d(s_1,s_2)+d(s_2,s_3)$. Extend our definition of Hamming distance to say that $d(\epsilon,\epsilon) = 0$, where $\epsilon \in \mathbb{Z}_2^0$ is the empty codeword (and in fact $\mathbb{Z}_2^0 = \{\epsilon\}$). We perform induction on the codeword length $n \in \mathbb{N}\cup \{0\}$. For our base case of $n=0$, we simply have that $d(\epsilon, \epsilon) \le d(\epsilon,\epsilon)+d(\epsilon, \epsilon)$, which is trivially true as $0 \le 0 + 0$. 
    
    Now we consider the inductive case for $n \ge 1$, and denote $s_1 = a_1:a_{2\dots n}$, $s_2 = b_1:b_{2\dots n}$ and $s_2 = c_1:c_{2\dots n}$, where $a_1,b_1,c_1 \in \mathbb{Z}_2$ are the first binary digits in each string and $a_{2\dots n},b_{2\dots n},c_{2\dots n} \in \mathbb{Z}_2^{n-1}$ are the rest of the digits in each codeword. By the definition of Hamming distance $d$, our statement $d(s_1,s_3)\le d(s_1,s_2)+d(s_2,s_3)$ can be written as \[|a_1 - c_1| + d(a_{2\dots n}, c_{2\dots n}) \le |a_1 - b_1| + d(a_{2\dots n},b_{2 \dots n}) + |b_1 - c_1| +d(b_{2\dots n},c_{2\dots n}).\]
    By the inductive hypothesis, we have that $d(a_{2\dots n}, c_{2\dots n}) \le d(a_{2\dots n},b_{2 \dots n}) + d(b_{2\dots n},c_{2\dots n})$, so we only need to show that $|a_1 - c_1|  \le |a_1 - b_1| +|b_1 - c_1|$. However, we can rewrite $|a_1 - c_1 | = |(a_1 - b_1) + (b_1 - c_1)|$, so we only need to show
    \[|(a_1 - b_1) + (b_1 - c_1)| \le |a_1 - b_1| +|b_1 - c_1|,\]
    but this is simply the regular triangle inequality $|x + y| \le |x| + |y|$. In our case, $x,y\in \{-1, 0, 1\}$ and we can do the casework. In fact, because in our case $x + y \in \{-1,0,1\}$ too, we do not need to consider the cases of $x = y = 1$ and $x = y = -1$. By the symmetry of the equations, WLOG, we do not need to consider permutations of $x$ and $y$. For $x = 0$ and $y = -1$, we get $|-1| \le |0| + |-1|$, for $x = 0$ and $y = 1$, we get $|1| \le |0| + |1|$, and finally for $x = -1$ and $y = 1$, we get $|0| \le |-1| + |1|$. QED.
    
\end{enumerate}
\end{Solution}

\vspace{0.5cm}

\begin{Solution}
{\bf (\S3.3, Problem 7)}\\
Let $f : A \to B$ and $g : B \to C$. 
\begin{enumerate}[(a)]
    \item If $g \circ f$ is injective, then $f$ is injective. Assume for contradiction that $f$ is not injective, so let $x,y \in A$ be elements such that $f(x)=f(y)$ and $x \ne y$. By the definition of $g \circ f$, we will have that $(g\circ f)(x)= g(f(x)) = g(f(y)) = (g\circ f)(y)$. Thus, we have a contradiction with the fact that $g \circ f$ is injective, as $x \ne y$ but $(g\circ f)(x) = (g\circ f)(y)$. Thus, our assumption was false and $f$ is injective.
    \item If $g \circ f$ is surjective, then $g$ is surjective. By the definition of surjectivity, this means that $\forall (c \in C), \exists(a \in A), g(f(a))=c$. We can directly use this to show that $g$ is surjective: let $c \in C$ be any element of $C$, and then we know that there is an $a\in A$ such that $f(a) \in B$ satisfies $g(f(a))=c$. Thus, we have shown that $\forall (c \in C), \exists(b \in B), g(b)=a$, and in particular $b = f(a)$ for some $a \in A$. We have shown that $g$ is surjective.
\end{enumerate}
\end{Solution}

\vspace{0.5cm}

\begin{Solution}
{\bf (\S3.3, Problem 8)}\\
\begin{enumerate}[(a)]
    \item Consider $f : \mathbb{R} \to \mathbb{R}$ and $g : \mathbb{R} \to \mathbb{R}$ given by $f(x) = e^x$ and $g(x)=x^2$. We see that $g$ is clearly not injective because $g(-1)=g(1)=1$, but $g \circ f$ is injective as $(g\circ f)(x) = e^{2x}$, and in fact it is a bijection between $\mathbb{R}$ and $\mathbb{R}_{> 0}$, with the inverse given by $\frac{1}{2} \ln x$.
    \item Consider $f : \mathbb{R} \to \mathbb{R}$ and $g : \mathbb{R} \to \{ 0\}$ given by $f(x) = \sin(x)$ and $g(x) = 0$. We see that $f$ is not surjective because for $y=10$, there does not exist any $x \in \mathbb R$ such that $f(x)=y$ as $\sin(x)\in [-1,1]$. However, $g\circ f : \mathbb{R} \to \{0\}$ clearly is surjective: we only need to consider $y=0$, and $x=0$ satisfies $(g\circ f)(x)=0$ (in fact, all $x \in \mathbb R$ do!).
\end{enumerate}
\end{Solution}

\vspace{0.5cm}

\begin{Solution}
{\bf (\S3.3, Problem 10)}\\
\begin{enumerate}[(a)]
    \item Epimorphisms: Let $f: A \to B$ be surjective but not injective. The left cancellation law does not hold: $f \circ g = f \circ h \implies g = h$. Consider $f$ to be a surjective constant function $f : A \to \{0\}$. Then $f \circ g = f \circ h$ will hold for any $g, h : C \to A$ whatsoever, as $\forall (x \in C), f(g(x)) = f(h(x))=0$. We can have $g \ne h$ and this still works!
    
    However, the right cancellation law still holds: $r \circ f = s \circ f \implies r = s$. We have that $r$ and $s$ must have domain $B$ in order for this composition to work,\footnote{This is actually a really weird (but valid) move: since we define functions to be equal to their graphs, using the codomain is actually using some data that is external to the functions (and both function composition and the concept of surjectivity inherently rely on this extra data!). We need to make this move to prevent $r$ and $s$ from possibly being defined on some random elements outside of $B$.} so we only need to show that for all $b \in B$, we have $r(b) = s(b)$. By the surjectivity of $f$, we can rewrite this $b$ as $b = f(a)$ for some $a \in A$, and we only need to show that $r(f(a))=s(f(a))$. However, this follows immediately from our hypothesis that $r \circ f = s \circ f$, so we have that $r = s$ and the right cancellation law still holds.

    \item Monomorphisms: Let $f : A \to B$ be injective but not surjective. The right cancellation law does not hold: $r \circ f = s \circ f \implies r = s$. Consider $f$ to be an inclusion map $f : [0,1] \xhookrightarrow{} \mathbb{R}$ and $r,s : \mathbb{R} \to \mathbb{R}$ to be $r(x)=x$ and $s(x)=|x|$. Clearly, $r\ne s$, but $r \circ f = s\circ f$ and in fact $r \circ f = s \circ f = f$.

    However, the left cancellation law does hold: $f \circ g = f \circ h \implies g = h$. Let $g$ and $h$ be functions $C \to A$. By our hypothesis, we have that for all $c \in C$, $f(g(c)) = f(h(c))$. By the injectivity of $f$, this equality implies that $g(c) = h(c)$, which immediately implies $g = h$.
\end{enumerate}
\end{Solution}

\vspace{0.5cm}

\begin{Solution}
{\bf (\S3.3, Problem 14)}\\
Let $f : A \to B$ and $g : C \to D$ be functions.
\begin{enumerate}[(a)]
    \item If $A \cap C = \varnothing$, then $(f \cup g) : A \cup C \to B \cup D$. We have that $f \subseteq A \times B$ and $g \subseteq C \times D$, so $f \cup g \subseteq (A \cup C) \times (B \cup D)$, and we have that $f \cup g$ is a relation on the desired sets. To show that it is a function, consider any $x \in A \cup C$. Because $A \cap C = \varnothing$, we have that either $x \in A$ or $x \in C$, but not both. If $x \in A$, then by the fact that $f$ is a function, we have a unique $y_b \in B$ such that $(x,y_b) \in f$, so we must have $(x,y_b) \in f \cup g$. Because $x \not\in C$, however, we must have that there does not exist any $y_d \in D$ such that $(x,y_d) \in g$, so $(x, y_b)$ must be the unique element with first element $x$ in $f \cup g$. Similar logic applies to the case where $x \in C$: we have a unique $y_d \in D$ such that $(x, y_d) \in g$, and since $x \not \in A$ there must not be any $y_b \in B$ such that $(x, y_b) \in f$, so $(x, y_d)$ is the unique element with first element $x$ in $f \cup g$. Thus, $f \cup g$ is a function.
    \item Consider the case where $A =B=\mathbb{R}$, and let $f(x)=x$ and $g(x)=-x$. Then we would have $(x,x) \in f \cup g$ and $(x, -x) \in f \cup g$, meaning $f \cup g$ is not a function because the $y$ in the codomain such that $(f\cup g)(x)=y$ is not unique.
    \item Let $f : (-\infty, 10] \to (-\infty, 10]$ be the identity on $(-\infty, 10]$ and similarly let $g : [-10, \infty) \to [-10,\infty)$ be the identity on $[-10,\infty)$. Then we have $A\cap C = (-\infty, 10] \cap [-10,\infty) = [-10,10] \ne \varnothing$, but we still have that $f \cup g : \mathbb R \to \mathbb R$ is a function, and in fact $f \cup g = \mathrm{id}_\mathbb{R}$. Sheaffification sheaffifies the sheaf.
    \item Assume that $f \cup g$ is a function. If $f$ is surjective and $g$ is surjective, then $f \cup g$ surjective. Let $y \in B \cup D$ be any element in $B \cup D$. By definition of set union, we either have that $y \in B$ or $y \in D$. Consider the $y \in B$ case. By the surjectivity of $f$, we have that there exists an $a \in A$ such that $f(a) = y$, so $(f\cup g)(a) = y$. In the case that $y \in D$, we have by the surjectivity of $g$ that there exists $c \in C$ such that $g(c) = y$, so $(f\cup g)(c) = y$. Thus, $f\cup g$ is surjective, as for all $y \in B \cup D$, there exists some $x \in A \cup C$ such that $(f\cup g)(x) = y$.
    \item If $f\cup g$ is an injective function, then $f$ and $g$ are injective. Assume FTSOC that there exist $a_1,a_2 \in A$ with $a_1 \ne a_2$ such that $f(a_1) = f(a_2)$. Then we must have that $(f\cup g)(a_1) = (f\cup g)(a_2)$, and by the injectivity of $f \cup g$, we get $a_1 = a_2$, a contradiction. Thus, our assumption was false and $f$ was injective. The same argument applies for $g$: assume FTSOC that $\exists(c_1,c_2 \in C), g(c_1) =g(c_2) \land c_1 \ne c_2$. Then $(f\cup g)(c_1) = (f\cup g)(c_2) \implies c_1 =c_2$, a contradiction. Thus, $f$ and $g$ are injective.
    \item If $f\cup g$ is a bijective function, then $(f\cup g)^{-1} = f^{-1} \cup g^{-1}$. Expanding definitions, we have \[(f \cup g)^{-1} = \{ (x,y) \mid (y,x) \in f\cup g \}\] and \[f^{-1} \cup g^{-1} = \{(x,y) \mid (y,x) \in f \} \cup \{(x, y)\mid (y,x) \in g \}.\] Thus, $(f\cup g)^{-1}$ and $f^{-1} \cup g^{-1}$ as equal as sets, as $(x,y) \in (f\cup g)^{-1}$ and $(x,y) \in f^{-1} \cup g^{-1}$ both amount to $(y,x) \in f \lor (y,x)\in g$ by the definition of set union.
    \item The example in part (c) works as an example of an injective $f \cup g$. Let $f = \mathrm{id}_{(-\infty,10]}$ and $g = \mathrm{id}_{[-10,\infty)}$. Then $f\cup g =\mathrm{id}_\mathbb{R}$, which is injective.
    \item Let $f : \mathbb{R}_{\le 0} \to \mathbb{R}$ be given by $f(x)=x$, and let $g : \mathbb{R}_{>0} \to \mathbb{R}$ be given by $g(x)=x-1$. Then $f$ and $g$ are both injective, and $f \cup g : \mathbb{R} \to \mathbb{R}$ is a function, but $f \cup g$ is not injective because $(f\cup g)(0) = (f\cup g)(1)$, but $0 \ne 1$. No gluability for presheaves...
\end{enumerate}
\end{Solution}


\end{document}
